\section*{Заключение}
\addcontentsline{toc}{section}{Заключение}

История энергоэффективных нейросетевых вычислений связана
с изменением представлений о затратах на обработку информации.
Математические модели нейронов и методы обучения определили,
какие преобразования должна выполнять сеть. Аналоговые машины
и ранние адаптивные схемы показали разные способы аппаратного
выполнения операций и хранения настраиваемых параметров.
Эти разработки решали различные задачи и не образовывали
единой последовательности нейросетевых устройств.

Работы Ландауэра и Беннетта позволили различать физические
ограничения при потере информации и затраты выбранного
способа вычислений. Комплементарная логика и масштабирование
Деннарда относились к другому вопросу: как уменьшать мощность
элементов и повышать их плотность при заданных условиях
работы схемы. Ограничения снижения напряжения усилили
значение организации вычислений и специализации аппаратуры.

Исследования быстрой локальной памяти, параллельного ускорения
и сложности обмена показали, почему характеристики отдельного
арифметического узла недостаточны для оценки всей задачи.
В нейроморфных схемах Мида и потоках данных Eyeriss локальность
связей и повторное использование значений учитывались
при выборе устройства вычислений. На этих примерах стоимость
передачи и хранения данных рассматривается совместно
со стоимостью арифметики.

Развитие сжатия, дистилляции и квантования связало сокращение
ресурсов с качеством решения прикладной задачи. Число весов,
число операций и разрядность стали параметрами выбора модели
и её реализации. Допустимость изменения этих параметров
проверяется по результатам работы сети, а энергетический
эффект — с учётом поддержки выбранного представления
памятью и исполнительными узлами.

Разработка методов сопровождалась уточнением правил
сравнения. Показатели энергии связывались с выполненной
работой, требованиями к качеству и условиями применения.
Моделирование и воспроизводимые эксперименты дополняли
описание новых алгоритмов и устройств. Поэтому история
энергоэффективных вычислений включает как развитие самих
методов, так и развитие способов оценки их результатов.
